\documentclass[conference, a4paper]{IEEEtran}

% --- PACKAGES REQUIS ---
\usepackage[utf8]{inputenc}
\usepackage[T1]{fontenc}
\usepackage[french]{babel}
\usepackage{amsmath, amssymb, amsfonts}
\usepackage{booktabs}
\usepackage{graphicx}
\usepackage{microtype}
\usepackage{cite}
\usepackage{hyperref}

% --- CONFIGURATION DU DOCUMENT ---
\title{Optimisation sous Contrainte d'Équité pour la Prédiction de Liens : Analyse Topologique et Robustesse sur Pokec}

% --- BLOC DES 4 AUTEURS ---
\author{
\IEEEauthorblockN{Younes Boutrouft\IEEEauthorrefmark{1},
Amir Kashani\IEEEauthorrefmark{1},
Mehdi Ait Hamma\IEEEauthorrefmark{7},
Tahiana Hajanirina Andriambahoaka\IEEEauthorrefmark{1}}
\IEEEauthorblockA{\textit{\IEEEauthorrefmark{1}MS Data \quad \IEEEauthorrefmark{7}MS IA}\\
\textit{Télécom Paris}}
}

\begin{document}

\maketitle

\begin{abstract}
Le réseau social Pokec présente un biais de genre marqué que les GCN reproduisent. On teste trois corrections sur la sous-région la plus biaisée : reweighting, régularisation DP, et FairGNN adversarial. Notre diagnostic XAI montre que retirer le feature genre \emph{aggrave} le biais : il vient de la topologie du graphe, pas du feature lui-même. La régularisation DP ($\lambda=0.1$) donne le meilleur compromis pratique. FairGNN ($\alpha=0.5$) va plus loin sur la fairness mais reste fragile à régler. MC Dropout et les tests de bruit confirment que les corrections survivent aux perturbations.
\end{abstract}

\begin{IEEEkeywords}
Équité Algorithmique, GCN, Prédiction de Liens, Apprentissage Antagoniste, Incertitude.
\end{IEEEkeywords}

\section{Introduction}
Les Graph Convolutional Networks \cite{kipf2017semi} sont devenus la référence pour la prédiction de liens dans les réseaux sociaux. Leur principe est simple : chaque nœud agrège les features de ses voisins via le message-passing. Sur un graphe homophile, ce mécanisme amplifie naturellement la similarité entre nœuds du même groupe. Notre question : ce biais structurel se traduit-il dans les prédictions, et que peut-on faire pour le corriger sans casser la performance ?

\section{Protocole Expérimental}
On évalue le pipeline sur \textbf{Pokec}, le plus grand réseau social slovaque, filtré sur la sous-région \textit{Ceska Republika} (31\,298 nœuds, 107\,802 arêtes, 64 features). C'est la région avec le biais de genre le plus marqué (ratio H$\to$F / H$\to$H = 4.59$\times$). L'attribut sensible est binaire (F=57\%, H=43\%).

\begin{table}[htbp]
\centering
\small
\begin{tabular}{ll}
\toprule
\textbf{Paramètre} & \textbf{Valeur} \\
\midrule
Encodeur & GCN 2 couches (64$\to$128$\to$64), Dropout 0.5 \\
Décodeur & Produit scalaire + sigmoïde \\
Split arêtes & Train 70\% / Val 15\% / Test 15\% \\
Optimiseur & Adam, weight\_decay=5e-4 \\
Learning rate & 0.01 (GCN, SGC) ; 0.005 (GAT) \\
Epochs & 100, best-model sur AUC val \\
Seed & 42 \\
Modèles comparés & GCN, SGC (K=2), GAT (4 têtes) \\
\bottomrule
\end{tabular}
\end{table}

\section{Diagnostic du Biais Topologique}

\subsection{Performance de Base et Fuite d'Information}
Le GCN baseline prédit bien (AUC = 0.855) mais avec un biais clair : DP gap = 0.051, juste au-dessus du seuil de fairness. En regardant les embeddings, les centroïdes des nœuds femmes et hommes pointent dans la même direction (cosine = 0.669). Le test le plus parlant : un simple classifieur linéaire entraîné sur les embeddings devine le genre à 82\% — bien au-dessus du baseline majoritaire à 57\%. C'est notre signal le plus net que le GCN encode le genre sans qu'on le lui ait demandé.

\subsection{Le Paradoxe de l'Ablation}
Pour isoler la source du biais, on a mis à zéro la colonne \texttt{gender} ($\mathbf{X}_{gender} = 0$) puis on a refait l'inférence. À notre surprise, le DP gap \emph{augmente} (0.051 $\to$ 0.063). Ce résultat est important : le biais ne vient pas du feature explicite, il vient de la topologie. Sans gender en entrée, le GCN reconstruit l'info via le message-passing sur un graphe homophile. Le GAT confirme cette lecture : son mécanisme d'attention pondère spontanément plus fort les voisins du même genre.

\section{Atténuation : Pénalité vs. Antagonisme}

Première leçon : le reweighting \cite{kamiran2012data} ne suffit pas. Le DP gap reste à 0.056 — pondérer la loss ne neutralise pas un biais qui vient de la structure du graphe. On se tourne donc vers deux approches \textit{In-processing}.

\subsection{Régularisation Continue de la Parité Démographique}
On ajoute une pénalité différentiable sur la différence des scores moyens entre paires intra et inter-genre \cite{agarwal2018reductions} :
\begin{equation}
\mathcal{L}_{total} = \mathcal{L}_{task} + \lambda \left| \mathbb{E}[\hat{p}_{intra}] - \mathbb{E}[\hat{p}_{inter}] \right|
\end{equation}
En faisant varier $\lambda$, on trace une courbe de Pareto. L'optimum est à $\lambda=0.1$ : DP gap divisé par deux (0.030) pour seulement 0.2\% de perte d'AUC (0.853). Au-delà de $\lambda = 1$, le modèle sur-corrige : le cosine bascule en négatif (les centroïdes pointent en directions opposées), et le Leakage reste élevé ($+30.3$ pp). On gagne en équité de sortie mais on perd en équité de représentation.

\subsection{Apprentissage Antagoniste (FairGNN)}
Pour nettoyer les embeddings du signal genre, FairGNN \cite{dai2021say} oppose l'encodeur ($E$) à un discriminateur ($A$) qui essaie de prédire le genre $s$ depuis les embeddings $\mathbf{Z}$. L'entraînement alterne un pas du discriminateur et un pas de l'encodeur :
\begin{equation}
\min_{\theta_E} \max_{\theta_A} \left( \mathcal{L}_{task} - \alpha \mathcal{L}_{adv}(A(\mathbf{Z}), s) \right)
\end{equation}
Avec $\alpha=0.5$, FairGNN donne le meilleur DP (0.004) pour une AUC de 0.848. Le cosine bascule en négatif, ce qui montre que les embeddings sont effectivement déstructurés. Mais le réglage de $\alpha$ est fragile : au-delà de 1, le DP remonte au-dessus du baseline. Ce comportement non-monotone est une limite à prendre au sérieux.

\begin{table}[htbp]
\centering
\caption{Synthèse des Performances et de la Robustesse}
\label{tab:results}
\begin{tabular}{lcccc}
\toprule
\textbf{Méthode} & \textbf{AUC} & \textbf{DP Gap} & \textbf{Leakage} & \textbf{DP (Bruit $\sigma=0.5$)} \\
\midrule
Baseline (GCN) & 0.855 & 0.051 & +25.0 pp & 0.010$^*$ \\
DP Reg ($\lambda=0.1$) & 0.853 & 0.030 & +30.3 pp & 0.040 \\
FairGNN ($\alpha=0.5$) & 0.848 & \textbf{0.004} & +34.8 pp & \textbf{0.015} \\
\bottomrule
\multicolumn{5}{l}{\scriptsize $^*$Le bruit détruit le signal topologique du baseline.}
\end{tabular}
\end{table}

\section{Incertitude et Robustesse}

\subsection{Quantification de l'Incertitude}
On utilise MC Dropout \cite{gal2016dropout} avec 30 passes (dropout actif à l'inférence). L'entropie prédictive moyenne est élevée ($\mathcal{H} \approx 0.56$, sur un maximum de 0.69), mais l'information mutuelle est faible ($\mathcal{I} \approx 0.025$). L'incertitude est donc surtout aléatoire — liée au bruit des données — pas épistémique. L'ECE vaut $\approx 0.16$ pour les trois modèles : le décodeur dot-product est mal calibré, indépendamment de la correction fairness.

\subsection{Résilience aux Perturbations}
On ajoute du bruit gaussien ($\sigma=0.5$) sur les features et on observe ce qui résiste. DP-reg ($\lambda=0.1$) voit son DP gap remonter à 0.040 — sa correction agit sur les scores de sortie, pas en profondeur. FairGNN ($\alpha=0.5$) reste sous 0.015. Il tient mieux parce qu'il a effectivement modifié les embeddings, pas juste recalibré les sorties.

\section{Utilisation d'outils d'IA}
Ce projet a utilisé \textbf{Claude} (Anthropic) comme assistant de développement pour la structuration du code, la recherche bibliographique et le debugging. Toutes les décisions méthodologiques et l'analyse des résultats sont le fruit du travail de l'équipe.

\section{Conclusion et Limites}
Notre analyse montre que le biais sur Pokec est avant tout topologique : retirer le feature genre l'aggrave, prouvant que le GCN le reconstruit via la structure du graphe. Côté correction, deux options se dégagent. DP-reg $\lambda=0.1$ est le bon compromis si on veut un modèle facile à régler et un coût en performance minimal. \textbf{FairGNN $\alpha=0.5$ offre le meilleur profil global} : DP gap quasi-nul (0.004), coût AUC modéré (-0.7\%), robustesse structurelle préservée. La contrepartie est un Cosine F/H négatif ($-0.16$), signe d'une représentation activement anti-homophilique — à surveiller pour les tâches downstream. Aucune méthode ne satisfait toutes les métriques en même temps, ce qui illustre le théorème d'impossibilité de Chouldechova \cite{chouldechova2017fair}.

\textbf{Limites.}
\begin{itemize}
\item Seed unique (42) : la variance inter-runs n'est pas quantifiée.
\item Attribut sensible binaire : pas d'intersectionnalité (genre $\times$ âge).
\item Calibration non corrigée : l'ECE $\approx 0.16$ pourrait être réduit par \textit{temperature scaling} \cite{guo2017calibration}.
\item Sensibilité au bruit gaussien sur les features : les 3 modèles passent sous le seuil de robustesse (Rob $\approx 0.94 < 0.95$) — vulnérabilité orthogonale au pipeline fairness.
\item FairGNN instable : sensibilité non-monotone à $\alpha$ (le DP gap augmente pour $\alpha \ge 2$).
\item Dataset Pokec (2012, Slovaquie) : généralisation non testée sur d'autres réseaux sociaux.
\end{itemize}

% --- BIBLIOGRAPHIE ---
\begin{thebibliography}{00}

\bibitem{kipf2017semi}
T.~N. Kipf and M.~Welling, ``Semi-supervised classification with graph convolutional networks,'' in \emph{International Conference on Learning Representations (ICLR)}, 2017.

\bibitem{dai2021say}
E.~Dai and S.~Wang, ``Say no to the discrimination: Learning fair graph neural networks with limited sensitive attribute information,'' in \emph{Proceedings of the 14th ACM International Conference on Web Search and Data Mining (WSDM)}, 2021, pp. 680--688.

\bibitem{gal2016dropout}
Y.~Gal and Z.~Ghahramani, ``Dropout as a bayesian approximation: Representing model uncertainty in deep learning,'' in \emph{International Conference on Machine Learning (ICML)}, 2016, pp. 1050--1059.

\bibitem{chouldechova2017fair}
A.~Chouldechova, ``Fair prediction with disparate impact: A study of bias in recidivism prediction instruments,'' \emph{Big Data}, vol.~5, no.~2, pp. 153--163, 2017.

\bibitem{kamiran2012data}
F.~Kamiran and T.~Calders, ``Data preprocessing techniques for classification without discrimination,'' \emph{Knowledge and Information Systems}, vol.~33, no.~1, pp. 1--33, 2012.

\bibitem{agarwal2018reductions}
A.~Agarwal, A.~Beygelzimer, M.~Dudik, J.~Langford, and H.~Wallach, ``A reductions approach to fair classification,'' in \emph{International Conference on Machine Learning (ICML)}, 2018, pp. 60--69.

\bibitem{guo2017calibration}
C.~Guo, G.~Pleiss, Y.~Sun, and K.~Q. Weinberger, ``On calibration of modern neural networks,'' in \emph{International Conference on Machine Learning (ICML)}, 2017, pp. 1321--1330.

\end{thebibliography}

\end{document}
